% `tex.web` §1054 `its_all_over`: `\end` does not end the job on the first try.
%
%     function its_all_over:boolean;
%     begin if privileged then
%       begin if (page_head=page_tail)and(head=tail)and(dead_cycles=0) then
%         begin its_all_over:=true; return; end;
%       back_input; {we will try to end again after ejecting residual material}
%       tail_append(new_null_box); width(tail):=hsize;
%       tail_append(new_glue(fill_glue));
%       tail_append(new_penalty(-'10000000000));
%       build_page;
%       end;
%     its_all_over:=false;
%     end;
%
% Unless the page and the current list are already empty, `\end` goes back on
% the input, the residual `\hbox to\hsize{} \vfill \penalty-'10000000000` is
% appended, and the page builder ejects the page — after which `main_control`
% reads `\end` a second time and finds nothing left. So `\tracingcommands`
% shows `{\end}` twice for any document that produced a page, and once for one
% that produced none.
%
% That second reading is also why nothing after `\end` in a file is ever read:
% the backed-up `\end` sits on top of the file, not after it.
%
% We used to clear the input stack outright and hand the residual material to a
% routine outside the main loop. The pages came out the same — which is why
% this only ever showed in the trace — but `\output` ran somewhere
% `main_control` was not.
\catcode`\{=1 \catcode`\}=2
\tracingcommands=1 \tracingonline=0
\font\r=cmr10 \r \hsize=100pt
A\par
\end
